\documentclass[11pt,a4paper]{article}
\newcommand{\intestazione}{Dimero di spin --- 10. VQE noise-aware}
% =====================================================================
%  Preambolo comune ai documenti sul dimero (serie "che cosa abbiamo fatto")
%  Incluso con % =====================================================================
%  Preambolo comune ai documenti sul dimero (serie "che cosa abbiamo fatto")
%  Incluso con % =====================================================================
%  Preambolo comune ai documenti sul dimero (serie "che cosa abbiamo fatto")
%  Incluso con \input{preambolo_dimero} subito dopo \documentclass.
% =====================================================================
\usepackage[T1]{fontenc}
\usepackage[utf8]{inputenc}
\usepackage[main=italian,provide=*]{babel}
\usepackage{amsmath,amssymb,mathtools}
\usepackage{graphicx}
\usepackage{booktabs}
\usepackage{colortbl}
\usepackage{xcolor}
\usepackage{geometry}
\usepackage{placeins}
\usepackage{caption}
\usepackage{enumitem}
\usepackage{fancyhdr}
\usepackage{needspace}
\usepackage{tikz}
\usetikzlibrary{arrows.meta,positioning,shapes.geometric,calc,fit,backgrounds}
\usepackage{hyperref}
\hypersetup{colorlinks=true,linkcolor=blu,citecolor=blu,urlcolor=blu}

\geometry{margin=2.5cm}

% --- colori coerenti con quelli delle figure (palette Okabe-Ito)
\definecolor{blu}{HTML}{0072B2}
\definecolor{arancio}{HTML}{E69F00}
\definecolor{verdes}{HTML}{009E73}
\definecolor{rossos}{HTML}{D55E00}
\definecolor{violas}{HTML}{CC79A7}
\definecolor{grigetto}{HTML}{E8E8E8}

\captionsetup{font=small,labelfont=bf,width=0.92\textwidth}

\pagestyle{fancy}
\fancyhf{}
\fancyhead[L]{\small\itshape\intestazione}
\fancyhead[R]{\small\thepage}
\renewcommand{\headrulewidth}{0.3pt}

% --- notazione
\newcommand{\ket}[1]{\lvert #1\rangle}
\newcommand{\bra}[1]{\langle #1 \rvert}
\newcommand{\media}[1]{\langle #1 \rangle}
\newcommand{\vs}[1]{\mathbf{s}_{#1}}

% --- riquadro "in una riga": il messaggio da portare a casa
\newcommand{\messaggio}[1]{%
  \begin{center}
  \begin{minipage}{0.93\textwidth}
  \setlength{\fboxsep}{9pt}%
  \colorbox{grigetto}{\begin{minipage}{\dimexpr\textwidth-2\fboxsep}
  \small\textbf{In una riga.}\enspace #1
  \end{minipage}}
  \end{minipage}
  \end{center}
}

% --- riquadro "come lo sappiamo": la verifica numerica dietro un'affermazione
\newcommand{\verifica}[1]{%
  \begin{center}
  \begin{minipage}{0.93\textwidth}
  \small\textcolor{verdes}{\rule{2.2em}{1.1pt}}\enspace
  \textcolor{verdes}{\textbf{Come lo sappiamo.}}\enspace #1
  \end{minipage}
  \end{center}
  \medskip
}
 subito dopo \documentclass.
% =====================================================================
\usepackage[T1]{fontenc}
\usepackage[utf8]{inputenc}
\usepackage[main=italian,provide=*]{babel}
\usepackage{amsmath,amssymb,mathtools}
\usepackage{graphicx}
\usepackage{booktabs}
\usepackage{colortbl}
\usepackage{xcolor}
\usepackage{geometry}
\usepackage{placeins}
\usepackage{caption}
\usepackage{enumitem}
\usepackage{fancyhdr}
\usepackage{needspace}
\usepackage{tikz}
\usetikzlibrary{arrows.meta,positioning,shapes.geometric,calc,fit,backgrounds}
\usepackage{hyperref}
\hypersetup{colorlinks=true,linkcolor=blu,citecolor=blu,urlcolor=blu}

\geometry{margin=2.5cm}

% --- colori coerenti con quelli delle figure (palette Okabe-Ito)
\definecolor{blu}{HTML}{0072B2}
\definecolor{arancio}{HTML}{E69F00}
\definecolor{verdes}{HTML}{009E73}
\definecolor{rossos}{HTML}{D55E00}
\definecolor{violas}{HTML}{CC79A7}
\definecolor{grigetto}{HTML}{E8E8E8}

\captionsetup{font=small,labelfont=bf,width=0.92\textwidth}

\pagestyle{fancy}
\fancyhf{}
\fancyhead[L]{\small\itshape\intestazione}
\fancyhead[R]{\small\thepage}
\renewcommand{\headrulewidth}{0.3pt}

% --- notazione
\newcommand{\ket}[1]{\lvert #1\rangle}
\newcommand{\bra}[1]{\langle #1 \rvert}
\newcommand{\media}[1]{\langle #1 \rangle}
\newcommand{\vs}[1]{\mathbf{s}_{#1}}

% --- riquadro "in una riga": il messaggio da portare a casa
\newcommand{\messaggio}[1]{%
  \begin{center}
  \begin{minipage}{0.93\textwidth}
  \setlength{\fboxsep}{9pt}%
  \colorbox{grigetto}{\begin{minipage}{\dimexpr\textwidth-2\fboxsep}
  \small\textbf{In una riga.}\enspace #1
  \end{minipage}}
  \end{minipage}
  \end{center}
}

% --- riquadro "come lo sappiamo": la verifica numerica dietro un'affermazione
\newcommand{\verifica}[1]{%
  \begin{center}
  \begin{minipage}{0.93\textwidth}
  \small\textcolor{verdes}{\rule{2.2em}{1.1pt}}\enspace
  \textcolor{verdes}{\textbf{Come lo sappiamo.}}\enspace #1
  \end{minipage}
  \end{center}
  \medskip
}
 subito dopo \documentclass.
% =====================================================================
\usepackage[T1]{fontenc}
\usepackage[utf8]{inputenc}
\usepackage[main=italian,provide=*]{babel}
\usepackage{amsmath,amssymb,mathtools}
\usepackage{graphicx}
\usepackage{booktabs}
\usepackage{colortbl}
\usepackage{xcolor}
\usepackage{geometry}
\usepackage{placeins}
\usepackage{caption}
\usepackage{enumitem}
\usepackage{fancyhdr}
\usepackage{needspace}
\usepackage{tikz}
\usetikzlibrary{arrows.meta,positioning,shapes.geometric,calc,fit,backgrounds}
\usepackage{hyperref}
\hypersetup{colorlinks=true,linkcolor=blu,citecolor=blu,urlcolor=blu}

\geometry{margin=2.5cm}

% --- colori coerenti con quelli delle figure (palette Okabe-Ito)
\definecolor{blu}{HTML}{0072B2}
\definecolor{arancio}{HTML}{E69F00}
\definecolor{verdes}{HTML}{009E73}
\definecolor{rossos}{HTML}{D55E00}
\definecolor{violas}{HTML}{CC79A7}
\definecolor{grigetto}{HTML}{E8E8E8}

\captionsetup{font=small,labelfont=bf,width=0.92\textwidth}

\pagestyle{fancy}
\fancyhf{}
\fancyhead[L]{\small\itshape\intestazione}
\fancyhead[R]{\small\thepage}
\renewcommand{\headrulewidth}{0.3pt}

% --- notazione
\newcommand{\ket}[1]{\lvert #1\rangle}
\newcommand{\bra}[1]{\langle #1 \rvert}
\newcommand{\media}[1]{\langle #1 \rangle}
\newcommand{\vs}[1]{\mathbf{s}_{#1}}

% --- riquadro "in una riga": il messaggio da portare a casa
\newcommand{\messaggio}[1]{%
  \begin{center}
  \begin{minipage}{0.93\textwidth}
  \setlength{\fboxsep}{9pt}%
  \colorbox{grigetto}{\begin{minipage}{\dimexpr\textwidth-2\fboxsep}
  \small\textbf{In una riga.}\enspace #1
  \end{minipage}}
  \end{minipage}
  \end{center}
}

% --- riquadro "come lo sappiamo": la verifica numerica dietro un'affermazione
\newcommand{\verifica}[1]{%
  \begin{center}
  \begin{minipage}{0.93\textwidth}
  \small\textcolor{verdes}{\rule{2.2em}{1.1pt}}\enspace
  \textcolor{verdes}{\textbf{Come lo sappiamo.}}\enspace #1
  \end{minipage}
  \end{center}
  \medskip
}

\newcommand{\Tr}{\operatorname{Tr}}

\title{\bfseries Vale la pena rioptimizzare vedendo il rumore?\\[2pt]
\large Documento 10 --- VQE noise-aware, estensione facoltativa}
\author{Samuele Schiavi \\ \small Tesi triennale in Fisica, Università di Parma
--- Relatore: Prof.\ Alessandro Chiesa}
\date{}

\begin{document}
\maketitle
\thispagestyle{fancy}

\section{La domanda complementare}

I cinque documenti precedenti (5--9) misurano quanto degrada un risultato
\textbf{ideale} --- parametri VQE ottimizzati senza rumore --- quando lo si
esegue su un simulatore rumoroso. È una scelta di modellazione dichiarata fin
dal Documento 6, non un'omissione: resta aperta la domanda complementare, se
il rumore fosse acceso già durante l'ottimizzazione --- come avverrebbe su un
hardware reale, dove non esiste una versione ideale a cui accedere per
calibrare l'ottimizzatore --- l'ottimizzatore convergerebbe a parametri
diversi? E se sì, questo cambierebbe in modo misurabile i risultati a valle
(Trotter, correlatori, lo scan del Documento 9)?

Questa è un'estensione \textbf{facoltativa} della pipeline, già chiusa in
tutti e cinque i documenti precedenti --- non uno stadio mancante. Copre: la
rioptimizzazione stessa, il suo effetto su $N^*$ della fedeltà di Trotter e su
$N^*$ del correlatore dinamico, la sua generalità rispetto al punto di lavoro
fisico e ai parametri di rumore.

\section{L'argomento fisico, e una previsione prima di verificarla}

Il canale di rumore è depolarizzante, isotropo: riduce uniformemente la sfera
di Bloch in ogni direzione. Per uno stato puro $\ket{\psi(\theta)}$ preparato
dall'ansatz,
\begin{equation}
\rho(\theta) = (1-\lambda)\,\ket{\psi(\theta)}\bra{\psi(\theta)} + \lambda\,\frac{\mathbf 1}{d},
\end{equation}
con $\lambda>0$ legato a $\varepsilon_{1q},\varepsilon_{2q}$ come nel
Documento 5, $d=4$ per 2 qubit. L'energia rumorosa diventa
\begin{equation}
E_\text{rumoroso}(\theta) = \Tr[\rho(\theta) H] = (1-\lambda)\,E_\text{ideale}(\theta) + \lambda\,\frac{\Tr[H]}{d}
\end{equation}
--- una trasformazione \textbf{affine} di $E_\text{ideale}(\theta)$: il secondo
termine non dipende da $\theta$, il primo è $E_\text{ideale}(\theta)$
moltiplicato per una costante positiva $1-\lambda$. Una trasformazione affine
con coefficiente positivo non cambia la posizione del minimo: lo stesso
$\theta^*$ che minimizza $E_\text{ideale}$ minimizza anche $E_\text{rumoroso}$.
A questo si aggiunge un secondo fatto, verificato indipendentemente
(Sez.~\ref{sec:cnot}): il conteggio CNOT del circuito transpilato non dipende
dai valori numerici dei parametri, quindi non c'è nemmeno una via indiretta,
via costo in gate, per cui convenga spostarsi verso un $\theta$ diverso da
quello ideale.

\messaggio{Previsione dichiarata prima della verifica: $\theta^*_\text{rumore}
\approx\theta^*_\text{ideale}$, nessun vantaggio misurabile a rioptimizzare
sotto rumore rispetto al semplice riuso dei parametri ideali. È una
previsione, non un fatto acquisito: il resto di questo documento la mette
alla prova su più metriche e più punti di lavoro.}

La previsione riguarda l'\emph{energia}. Ma l'effetto va controllato anche
sulle quantità a valle nella pipeline, ed è qui che serve una distinzione: la
fedeltà di Trotter (Documento 7) è un osservabile additivo,
$F(N)=\bra{\psi_\text{target}}\rho_N\ket{\psi_\text{target}}$; il correlatore
dinamico (Documento 8) usa invece $N^*=\arg\max_N|C(N)|$, il \textbf{modulo}
di un numero complesso. Sono oggetti matematici diversi --- un conto valido
per una somma additiva non si trasporta automaticamente a un modulo, come si
vedrà in modo netto nel Documento 11. Per questo entrambe le metriche vengono
verificate esplicitamente, non una per analogia dall'altra.

\section{Una scelta tecnica, scoperta lavorandoci}
\label{sec:cnot}

Il circuito va transpilato \textbf{dopo} aver assegnato i valori numerici di
$\theta$, ad ogni valutazione dell'obiettivo --- non transpilato una volta
sola in forma parametrica e poi solo legato con l'assegnazione dei parametri
(il principio inverso di quello che funziona per il blocco di Trotter, dove i
parametri fisici non cambiano da una valutazione all'altra). La ragione: con
parametri simbolici il transpilatore non può fondere le rotazioni a 1 qubit
il cui angolo dipende dal valore numerico di $\theta$, producendo un circuito
strutturalmente più rumoroso.

\begin{figure}[htbp]
\centering
\includegraphics[width=0.6\textwidth]{figure/fig_bug_conteggio_gate}
\caption{Costo in gate del circuito transpilato, a parità di parametri
$\theta$, a seconda dell'ordine fra transpilazione e assegnazione. Il
conteggio CNOT resta invariato (2); i gate a singolo qubit quasi raddoppiano
se si transpila prima dell'assegnazione --- 19 \texttt{rz} + 14 \texttt{sx}
contro 11 \texttt{rz} + 8 \texttt{sx} --- ogni gate in più è un'ulteriore
occasione di errore per il canale depolarizzante.}
\end{figure}

\FloatBarrier

\verifica{Invarianza del conteggio CNOT (premessa fisica dell'argomento
sopra) controllata su 5 assegnazioni casuali indipendenti di $\theta$: CNOT
$=2$ in tutte le prove, sempre 11 \texttt{rz} + 8 \texttt{sx} + 2 \texttt{cx}
--- nessuna direzione del paesaggio dei parametri verso cui convenga
spostarsi per ``risparmiare'' rumore.}

\section{Risultato 1: energia e fedeltà}

Multistart $R=6$ punti casuali più un seme uguale a
$\theta^*_\text{ideale}$, per garantire di non fare mai peggio del semplice
riuso; la robustezza del multistart è discussa a parte
(Sez.~\ref{sec:robustezza}).

\begin{center}
\renewcommand{\arraystretch}{1.25}
\begin{tabular}{@{}lcc@{}}
\toprule
Quantità & Riuso $\theta^*_\text{ideale}$ & VQE noise-aware\\
\midrule
$E$ (rumoroso) & $-3.50164605$ & $-3.50164616$\\
$F$ (rumorosa) & $0.98500601$ & $0.98500602$\\
CNOT & $2$ & $2$\\
\bottomrule
\end{tabular}
\end{center}

$\Delta E=-1.16\times10^{-7}$, $\Delta F=+1.76\times10^{-8}$ --- sei ordini di
grandezza più piccoli della scala fisica del problema (la perdita di fedeltà
$\sim1.5\times10^{-2}$ già misurata nel Documento 6 a rumore acceso).

\begin{figure}[htbp]
\centering
\includegraphics[width=0.92\textwidth]{figure/fig_scan_eps2q_confronto}
\caption{Energia e fedeltà rumorose in funzione di $\varepsilon_{2q}$, riuso
dei parametri ideali contro rioptimizzazione noise-aware. Le due curve sono
indistinguibili a questa scala su tutto il range esplorato, dal rumore nullo a
$4\times$ il valore di riferimento.}
\end{figure}

\FloatBarrier

\section{Robustezza del multistart}
\label{sec:robustezza}

Con $R=6$ punti casuali e nessun seme, il minimo trovato può restare a
distanza $10^{-4}$--$10^{-2}$ dall'ottimo vero, a seconda del seme ---
dispersione reale del paesaggio energetico rumoroso, caratterizzata su 8 semi
indipendenti. Per questo il risultato principale include
$\theta^*_\text{ideale}$ come uno dei punti di partenza.

\begin{figure}[htbp]
\centering
\includegraphics[width=0.72\textwidth]{figure/fig_robustezza_multistart}
\caption{Energie finali dei singoli restart COBYLA (rumore acceso), per tre
configurazioni di multistart. La linea tratteggiata è l'ottimo noto. Con soli
6 restart casuali il minimo trovato può restare lontano dall'ottimo; con 20 il
punto migliore si avvicina entro $3.6\times10^{-6}$.}
\end{figure}

\FloatBarrier

Un multistart indipendente più ampio ($R=20$, nessun seme) converge comunque
allo stesso ottimo entro $2.6\times10^{-6}$: il seme non altera la conclusione
fisica, garantisce solo di raggiungerla con un numero contenuto di restart.

\section{Risultato 2: $N^*$ della fedeltà di Trotter}

\begin{center}
\renewcommand{\arraystretch}{1.2}
\begin{tabular}{@{}lcc@{}}
\toprule
Preparazione & $N^*$ & $F(N^*)$\\
\midrule
$\theta^*_\text{ideale}$ & $8$ & $0.817725$\\
$\theta^*_\text{noise-aware}$ & $8$ & $0.817767$\\
\bottomrule
\end{tabular}
\end{center}

\begin{figure}[htbp]
\centering
\includegraphics[width=0.92\textwidth]{figure/fig_F_vs_N_confronto}
\caption{Sinistra: fedeltà $F(N)$, le due preparazioni a confronto ---
massimo a $N^*=8$ in entrambi i casi. Destra: scostamento, scala ingrandita
($\sim10^{-4}$): un ordine di grandezza sopra il rumore numerico, ma sei
ordini sotto la scala fisica del problema.}
\end{figure}

\FloatBarrier

\section{Risultato 3: $N^*$ del correlatore dinamico}

Per la distinzione strutturale già discussa, questa metrica va verificata
\textbf{direttamente}, non per estensione dal risultato precedente.

\begin{center}
\renewcommand{\arraystretch}{1.2}
\begin{tabular}{@{}lcc@{}}
\toprule
Preparazione & $N^*$ & $|C(N^*)|$\\
\midrule
$\theta^*_\text{ideale}$ & $5$ & $0.250002$\\
$\theta^*_\text{noise-aware}$ & $5$ & $0.249933$\\
\bottomrule
\end{tabular}
\end{center}

\begin{figure}[htbp]
\centering
\includegraphics[width=0.92\textwidth]{figure/fig_correlatore_noise_aware_confronto}
\caption{Sinistra: modulo del correlatore rumoroso, le due preparazioni a
confronto --- massimo a $N^*=5$ in entrambi i casi. Destra: scostamento,
scala ingrandita ($\sim10^{-5}$) --- un ordine di grandezza più piccolo
persino dello scostamento osservato per la fedeltà di Trotter.}
\end{figure}

\FloatBarrier

\section{Perché il risultato generalizza: l'argomento di continuità}

Le due metriche condividono una proprietà comune, verificabile direttamente:
$\theta^*_\text{noise-aware}$ e $\theta^*_\text{ideale}$ sono quasi
indistinguibili ($\Delta E\sim10^{-7}$). Per qualunque funzione ragionevolmente
liscia dello stato preparato --- una fedeltà, il modulo di un correlatore, o
altro --- questo garantisce, per continuità, uno scostamento a valle
altrettanto minuscolo, indipendentemente dalla struttura specifica della
metrica. Non è quindi la struttura di $N^*$ a garantire l'invarianza: è la
piccolezza di $\Delta\theta$ a renderla ininfluente su qualunque metrica
ragionevole a valle. Questo argomento è stato reso quantitativo: lo
scostamento osservato in entrambe le metriche è riprodotto entro lo $0.1\%$
da una previsione lineare al prim'ordine,
$\Delta f\approx\nabla f(\theta^*_\text{ideale})\cdot\Delta\theta$ (derivazione
completa in una nota di lavoro interna, non riportata per esteso qui per la
sua natura metodologica più che di risultato fisico).

\section{Generalità: un secondo punto di lavoro}

L'argomento fisico sopra è generale --- vale per qualunque $b/J,D/J$ e
qualunque intensità di rumore isotropo --- ma la \emph{dimensione} di
$\Delta\theta$, e quindi se sia abbastanza piccola da non spostare un $N^*$
specifico, è un fatto numerico. Per non lasciarlo verificato in un solo punto,
l'intera catena --- VQE ideale, VQE noise-aware, $N^*$ Trotter, $N^*$
correlatore --- è stata ripetuta a un secondo punto di lavoro, $b/J=-0.18$,
$D/J=1$: un termine DM cinque volte più forte e segno di $b$ opposto rispetto
a ``test 2''.

\begin{center}
\footnotesize
\renewcommand{\arraystretch}{1.25}
\begin{tabular}{@{}lcc@{}}
\toprule
Quantità & test 2 & secondo punto\\
 & ($b/J{=}0.35,D/J{=}0.80$) & ($b/J{=}{-}0.18,D/J{=}1$)\\
\midrule
$\Delta E$ (noise-aware $-$ riuso ideale) & $-1.16\times10^{-7}$ & $-2.56\times10^{-8}$\\
$\Delta F$ (noise-aware $-$ riuso ideale) & $+1.76\times10^{-8}$ & $+3.96\times10^{-9}$\\
$N^*$ Trotter (ideale $\to$ noise-aware) & $8\to8$ & $8\to8$\\
$N^*$ correlatore (ideale $\to$ noise-aware) & $5\to5$ & $3\to3$\\
\bottomrule
\end{tabular}
\end{center}

\begin{figure}[htbp]
\centering
\includegraphics[width=0.92\textwidth]{figure/fig_secondo_punto_confronto}
\caption{Secondo punto di lavoro. Sinistra: fedeltà di Trotter, $N^*=8$ ---
stesso valore del punto ``test 2''. Destra: modulo del correlatore, $N^*=3$
--- diverso dal punto ``test 2'' ($N^*$ dipende dalla fisica del punto, non è
universale), ma comunque invariato fra le due preparazioni.}
\end{figure}

\FloatBarrier

$N^*$ \textbf{cambia} da un punto di lavoro all'altro: è una proprietà della
fisica del punto, non della preparazione VQE. Quello che \textbf{non} cambia,
in entrambi i punti testati, è che $N^*$ resta lo stesso confrontando le due
preparazioni --- coerente con l'argomento di continuità: $\Delta\theta$ resta
piccolo anche a questo secondo punto, pur essendo un regime fisicamente
diverso. Due punti non dimostrano una legge generale per ogni $b/J,D/J$, ma
escludono che il risultato sia un caso isolato.

\section{Generalità: tutta la griglia dello scan sui parametri di rumore}

Lo scan sui parametri di rumore (Documento 9) esplora una griglia
$(\varepsilon_{1q},\varepsilon_{2q},p_\text{readout})$. Le due metriche a
valle sono state ricalcolate con \textbf{entrambe} le preparazioni su
\textbf{tutta} la griglia già usata in quello scan --- non solo al punto di
riferimento --- per un totale di 25 combinazioni.

\begin{figure}[htbp]
\centering
\includegraphics[width=0.95\textwidth]{figure/fig_passo5_griglia_completa}
\caption{$N^*$ in funzione di $\varepsilon_{2q}$, riuso ideale contro
noise-aware, su tutta la griglia del Documento 9. Sinistra: fedeltà di
Trotter (con la nota transizione a $N^*=1$ per rumore forte, già
caratterizzata lì). Destra: modulo del correlatore. Le due curve sono
sovrapposte in ogni punto.}
\end{figure}

\FloatBarrier

$N^*$ resta invariato fra le due preparazioni in \textbf{tutti} i 25 punti
testati --- 7 valori di $\varepsilon_{2q}$ e 6 di $\varepsilon_{1q}$ per la
fedeltà di Trotter, 5 valori di $\varepsilon_{2q}$, 4 di $\varepsilon_{1q}$ e 3
di $p_\text{readout}$ per il correlatore --- verificato punto per punto, non
solo al riferimento.

\section{Considerazioni}

\begin{itemize}[leftmargin=1.5em,itemsep=3pt]
\item La previsione è confermata, non solo plausibile. L'argomento della
trasformazione affine non garantiva a priori uno scostamento numerico
trascurabile su \emph{ogni} metrica a valle: il modulo di un correlatore e
una fedeltà additiva sono oggetti diversi, e comporre una trasformazione con
un modulo non eredita automaticamente le stesse proprietà di invarianza (si
vedrà in modo netto nel Documento 11). Per questo entrambe le metriche sono
state verificate esplicitamente, non una per analogia dall'altra.
\item L'argomento di continuità è più solido di un argomento di struttura, e
più generale: si applica a qualunque osservabile liscio dello stato preparato,
additivo o no, purché il gradiente resti moderato --- non richiede che
l'osservabile abbia una decomposizione additiva nota a priori.
\item Il dominio di validità è dichiarato esplicitamente: due punti di
lavoro, un insieme di parametri di rumore esplorato per intero al punto di
riferimento del canale. Non è una dimostrazione per ogni $b/J,D/J$ o ogni
intensità di rumore, ma un controllo sistematico, non un singolo punto
isolato.
\end{itemize}

\messaggio{Rioptimizzare l'ansatz PMA-2q$\cdot$3 sotto rumore, invece di
riusare i parametri ideali, non cambia in modo misurabile né l'energia o la
fedeltà del Documento 6 ($\Delta E\sim10^{-7}$--$10^{-8}$ a seconda del punto
di lavoro), né la posizione di $N^*$ della fedeltà di Trotter (Documento 7,
sempre $N^*=8$) o del correlatore dinamico (Documento 8, $N^*=5$ o $N^*=3$ a
seconda del punto di lavoro, comunque invariato). Verificato su due metriche
strutturalmente diverse, due punti di lavoro fisici, e l'intera griglia di
parametri di rumore del Documento 9 (25/25 punti) --- non un'estrapolazione da
un singolo controllo.}

\end{document}
